Дана строка с путём "/home/user/docs" (Unix). Добавить "/" в конец, если его нет.
Дана строка с путём "/home/user/docs/" (Unix). Убрать "/" из конца, если он есть.
Дана строка с путём "/home/user/docs/" (Unix). Подсчитать количество каталогов в пути.
Дана строка с путём "/opt/cprocsp/bin/amd64/" (Unix). Вывести имя самого верхнего каталога (opt).
Дана строка с путём "/opt/cprocsp/bin/amd64/" (Unix). Вывести имя самого нижнего каталога (amd64).
Дана строка с путём "C:\\Users\\file.txt" (Windows). Заменить "\\" на "/".
Дана строка с путём "/home/user/docs/file.txt" (Unix). Заменить "/" на "\\".
Дана строка с путём "/home/user/docs/file.txt" (Unix). Вывести имя файла (file.txt).
Дана строка с путём "/home/user/docs/file.txt" (Unix). Вывести расширение файла (txt).
Дана строка с путём "/home/user/docs/file.txt" (Unix). Вывести путь без имени файла (/home/user/docs/).
Дана строка с путём "/home/user/docs/file.txt" (Unix). Вывести имя файла без расширения (file).
Дана строка с путём "file.txt" (Unix). Определить, является ли путь абсолютным.
Дана строка с путём "/home/user/docs" (Unix). Определить, является ли путь относительным.
Дана строка с путём "C:\\Users\\docs\\file.txt" (Windows). Вывести имя диска (C:).
Дана строка с путём "/home/user/docs/" (Unix). Проверить, заканчивается ли путь на "/".
Дана строка с путём "/home//user///docs/" (Unix). Нормализовать путь, убрав дублирующиеся слэши (/home/user/docs/).
Дана строка с путём "/home/user/docs/" (Unix). Вывести путь к родительскому каталогу (/home/user/).
Дана строка с путём "/home/user/docs/file.txt" (Unix). Заменить имя файла на "report.pdf" (/home/user/docs/report.pdf).
Дана строка с путём "C:\\Program Files\\App\\app.exe" (Windows). Вывести имя самого нижнего каталога (App).
Дана строка с путём "/var/www/html/index.html" (Unix). Проверить, имеет ли файл расширение "html".
Дана строка с путём "/home/user/docs/" (Unix). Вывести все каталоги пути через запятую (home, user, docs).
Дана строка с путём "C:\\Users\\file.txt" (Windows). Определить, является ли путь абсолютным (Windows).
Дана строка с путём "docs/file.txt" (Unix). Преобразовать относительный путь в абсолютный, добавив "/home/user/" в начало.
Дана строка с путём "/home/user/docs/file.txt" (Unix). Проверить, содержит ли путь каталог "docs".
Дана строка с путём "/home/user/docs/file.txt" (Unix). Вывести глубину вложенности файла (количество каталогов).
Дана строка с путём "C:\\Users\\docs\\file.txt" (Windows). Вывести путь без имени диска (\\Users\\docs\\file.txt).
Дана строка с путём "/home/user/docs/" (Unix). Добавить "file.txt" в конец пути.
Дана строка с путём "/home/user/docs/" (Unix). Проверить, пуст ли путь (не содержит каталогов и файлов).
Дана строка с путём "/home/user/docs/" (Unix). Вывести предпоследний каталог в пути (user).
Дана строка с путём "/home/user/file.tar.gz" (Unix). Вывести последнее расширение файла (gz).
